\begin{abstract}
Self-supervised learning (SSL) faces a fundamental conflict between semantic understanding and image reconstruction. High-level semantic SSL (e.g., DINO) relies on global tokens that are forced to be location-invariant for augmentation alignment, a process that inherently discards the spatial coordinates required for reconstruction. Conversely, generative SSL (e.g., MAE) preserves dense feature grids for reconstruction but fails to produce high-level abstractions. We introduce STELLAR, a framework that resolves this tension by factorizing visual features into a low-rank product of semantic concepts and their spatial distributions. This disentanglement allows us to perform DINO-style augmentation alignment on the semantic tokens while maintaining the precise spatial mapping in the localization matrix necessary for pixel-level reconstruction. We demonstrate that as few as 16 sparse tokens under this factorized form are sufficient to simultaneously support high-quality reconstruction (2.60 FID) and match the semantic performance of dense backbones (79.10\% ImageNet accuracy). Our results highlight STELLAR as a versatile sparse representation that bridges the gap between discriminative and generative vision by strategically separating semantic identity from spatial geometry. Code available at \url{https://aka.ms/stellar}.

% Self-supervised vision encoders have become critical components of modern machine learning systems. Despite remarkable advances in image understanding, generation, and multimodal alignment, the underlying representation of visual features has remained largely unchanged, constrained by historical architectures and benchmarks. This reliance on dense feature grids introduces redundancy and limits the integration of understanding and generation. We propose a novel framework that represents images with a small number of sparse tokens \rev{in the form of} low-rank matrix factorization. While mathematically simple, this formulation effectively disentangles semantic and spatial information. We demonstrate that vision-only self-supervised learning under this framework yields sparse token representations that simultaneously support high-quality image understanding, detailed pixel-level reconstruction, and fine-grained semantic understanding. Together, these results highlight sparse tokens as a promising alternative to dense grids for efficient and versatile visual representation learning.

\end{abstract}